\section{A counterexample to McNew's natural-density conjecture}\label{sec:rankin}
For $E\subseteq\mathbb N$, its natural density is $d(E)=\lim_{x\to\infty}|E\cap[1,x]|/x$ when the limit exists. Let $\alpha$ be the supremum of these densities over sets avoiding nonconstant rational-ratio three-term progressions. McNew conjectured that Rankin's set attains $\alpha$ \cite[Conj.~7.2]{mcnew}. Let $G$ be the nonnegative integers whose ternary digits are all $0$ or
$1$. Rankin's set is $R=\{n\ge1:v_p(n)\in G\text{ for every prime }p\}$.
The set $G$ has no nonconstant three-term arithmetic progression:
in $u+w=2v$, the lowest ternary digit at which the three values differ
would give a nonconstant solution with digits in $\{0,1\}$ modulo $3$.
The valuations of a rational-ratio geometric progression satisfy this
identity at every prime. Thus $R$ is progression-free.
It has natural density approximately $0.71974$. Riddell had improved the corresponding constructions for lengths at least four \cite{riddell}. Here the length is three.

\medskip\noindent\textbf{Theorem.}
There is a subset $A$ of the positive integers with a natural density
strictly greater than $d(R)$ and with no nonconstant three-term
geometric progression of rational ratio. In particular
$\alpha>d(R)$, which disproves Conjecture 7.2 of McNew's published paper.

\medskip\noindent\textit{Construction.}
Choose a finite set $B$ of odd primes such that
$\sum_{q\in B}1/q>13$. One definite choice is the shortest initial
segment of the odd primes with this property. Euler's divergence theorem
ensures that the segment is finite. Alternatively use every odd prime up
to the explicit integer $Y_0=\lceil\exp(\exp(15))\rceil$. The finite Euler
product over $p\le Y_0$ is at least $\sum_{n\le Y_0}1/n\ge\log Y_0$. In its
logarithm the prime powers of exponent at least two contribute at most
$\sum_{n\ge2}1/[n(n-1)]=1$. Thus the odd-prime reciprocal sum up to
$Y_0$ is at least $\log\log Y_0-1-1/2\ge13.5$. Put $P=2\prod_{q\in B}q$ and
$S=\{2\}\cup B$. For the moment work solely with the semigroup of
$S$-smooth positive integers. Write $C$ for Rankin's smooth integers.
Make the following deletions from $C$.
\begin{itemize}
\item Delete those whose exponent equals $1$ at every prime of $B$.
\item For each $q\in B$, delete those whose exponent is at least $3$ at
$q$ and equals $1$ at every other prime of $B$.
\end{itemize}
In both rules the exponent at $2$ can be any member of $G$.
To the survivors add $X=\{2Pq:q\in B\}$ and call the resulting smooth
set $F$. Each added integer has exponent $2$ at $2$ and at exactly one
odd prime $q$, and exponent $1$ at all other odd primes in $B$.

\medskip\noindent\textit{Progression avoidance.}
Three retained old points were already in $C$. Suppose a progression
contained exactly one new point, $2Pq$. At every $r\in B\setminus\{q\}$
the three exponents belong to $G$ and one equals $1$, so all three equal
$1$. In this slice the deletion rules leave an old point with exponent
at $q$ only if that exponent is $0$: exponent $1$ was deleted by the
first rule, $2\notin G$, and exponents at least $3$ were deleted by the
second. The three $q$-exponents would therefore be $0,0,2$ in some
order, which cannot be an arithmetic progression. With two new points,
two $2$-adic exponents equal $2$, forcing the third to equal $2$;
an old point cannot do so. Three new points can form no nonconstant
progression because every odd coordinate lies in $\{1,2\}$ and the
$2$-adic coordinate is constant. These cases permit any ordering of the
three terms and therefore cover all rational ratios.

\medskip\noindent\textit{Gain in reciprocal mass.}
For a set of $S$-smooth integers put $W(E)=\sum_{n\in E}1/n$;
the sum converges because $S$ is finite.
The added mass is $P^{-1}\sum_{q\in B}1/(2q)$.
For deleted points of the first type the odd part is $P/2$,
so their mass in units of $1/P$ is
$2\sum_{e\in G}2^{-e}\le4$. For the second type, raising the
$q$-exponent from $1$ to $e\ge3$ multiplies the reciprocal weight
by $q^{1-e}$. Summing over all $e\ge3$ gives $1/[q(q-1)]$.
It follows that
\[
 P W(C\setminus F)
 \le 4\left(1+\sum_{q\in B}\frac{1}{q(q-1)}\right)
 \le 4\left(1+\sum_{n=3}^\infty\frac{1}{n(n-1)}\right)=6.
\]
No new point belongs to $C$. Consequently
\[
                 W(F)-W(C)>\frac{1}{2P}.
\]
The saving comes from paying once for the deleted slice with all odd
exponents equal to $1$. Each new prime contributes $1/(2Pq)$, while
the additional deletion cost at that prime is bounded by
$4/[Pq(q-1)]$. Requiring every addition to have $2$-adic exponent
exactly $2$ controls progressions involving two additions simultaneously.

\medskip\noindent\textit{Natural density.}
Let $H$ be the positive integers coprime to $P$ with $v_p(h)\in G$
at every prime outside $S$. They have positive natural density
\[
 c_H=\prod_{p\in S}(1-p^{-1})
      \prod_{p\notin S}\left((1-p^{-1})\sum_{e\in G}p^{-e}\right),
\]
where the second product multiplies the whole local factor
$(1-p^{-1})\sum_{e\in G}p^{-e}$ at each $p\notin S$.
Each factor in the second product is at least $1-p^{-2}$, since $0,1\in G$.
Finite-prime valuation frequencies follow from the Chinese remainder
theorem by truncating the exponents. Omitting conditions at primes
$p>Y$ costs upper density at most $\sum_{p>Y}p^{-2}\to0$, proving the
product formula. Set $A=\{hf:h\in H,\ f\in F\}$. In any progression
the valuations outside $S$ are constant by the progression-freeness of
$G$, so the three rough parts $h$ agree. Cancelling $h$ reduces the
progression to one in $F$. This proves the global avoidance claim.

By unique factorization the sets $fH$ are disjoint for distinct smooth
$f$, and $d(fH)=c_H/f$. For each $K$ the union over $f>K$ contributes
upper density at most $\sum_{f>K}1/f$, tending to zero. Hence $A$ has
natural density $c_H W(F)$, while $d(R)=c_HW(C)$ by the identical
argument. In particular
\[
                         d(A)-d(R)>\frac{c_H}{2P}>0 .
\]
For the explicit cutoff just given, $c_H\ge(\varphi(P)/P)/\zeta(2)>1/(2P)$
and $P\le Y_0!$, so the gain even exceeds $1/[4(Y_0!)^2]$.
The bound proves strict improvement; its deliberately large finite
prime set makes it unsuitable for a visible change in the decimal
approximation $0.71974$. The greatest possible natural density remains
to be determined. This proof uses Euler's theorem and elementary
valuation arguments; no computational search or zero theorem enters it.

\paragraph{Why optimizing a fixed finite list of primes can miss the improvement.}
Fix a finite set of primes $T$, a nonempty progression-free family $D$ of $T$-smooth integers, and Rankin's original rule at all primes outside $T$. The resulting product construction always admits a strict improvement. Indeed, choose a marker prime $\ell\notin T$ and put all satellites $q\in B$ outside $T\cup\{\ell\}$. Replacing $2$ by $\ell$ in the proof makes the normalized added mass $\ell^{-1}\sum_{q\in B}q^{-1}$ and bounds the deleted mass by $\ell(1-1/\ell)^{-1}(1+\sum_{q\in B}1/[q(q-1)])$. The first quantity is unbounded as $B$ grows. The two unchanged factors, namely $D$ and the remaining unmodified primes, multiply the gain by a positive constant in the density formula. A progression would project to a progression in each block, forcing constancy in every block. All changed integers can also be required to exceed any prescribed cutoff by choosing $B$ beyond it. This explains how exact finite-prime optimality results can coexist with failure of the unrestricted conjecture.

\paragraph{Relation to the previous approaches.}
Rankin's independent prime-exponent rule, McNew's comparisons between ratio conventions, and the prior report's sharp four-prime result all use the same valuation identity. The Lean proof's sparse-place argument is another instance of converting multiplicative patterns to forbidden additive relations. The step here is to \emph{correlate} prime exponents: fix one forbidden marker and share a single deletion among additions indexed by many primes. No estimate for zeros or theta lifts is required.
